% === VI. EXPERIMENTS AND RESULTS =============================================
\section{Experiments and Results}
\label{sec:results}

\subsection{Setup}
\label{subsec:setup}

All experiments run on the architecture of Section~\ref{sec:platform}: the
Formula AI vehicle in CARLA~\cite{carla2017} driven through ROS~2 Humble, with
the server in synchronous mode at a \SI{0.033}{\second} fixed physics step. The
identifier collects \SI{30}{\second} of data at \SI{50}{\hertz} and
re-identifies from the newest buffer, matching the baseline's data
budget~\cite{ontrack2025}. Ground truth for forces, slip angles and normal
loads is the simulator's own per-wheel telemetry, which the identifier never
reads.

\subsubsection{Surface}
Every run reported here uses a \emph{static} friction surface: the per-wheel
coefficient is held at its configured value of \num{1.0} for the whole run, and
the road-surface factor of \num{0.70} that PhysX applies on top of it makes the
effective coefficient the identifier is scored against \num{0.70}, constant to
the last digit the simulator reports over all \num{10744} and \num{12025}
force-validation samples of the two runs. Time-varying friction is out of scope
here (Section~\ref{sec:conclusion}).

This surface is deliberately \emph{not} the one the tire prior was measured on.
The prior of Section~\ref{subsubsec:priorprovenance} was fitted on a
$\mu=\num{1.05}$ surface and carries $D_{f}=\num{1.009}$,
$D_{r}=\num{1.002}$, so it overstates this run's grip by \SI{44}{\percent},
and it understates the axle cornering stiffness the plant realises here by
\SI{24}{\percent} at the front and \SI{16}{\percent} at the rear. A pipeline that
cannot move off its prior therefore has nowhere to hide here, which is the
point of running on a surface the prior does not describe. It remains a
\emph{single} surface, and the adaptivity claim that a friction sweep would
support is not made (Section~\ref{subsec:limits}).

\subsubsection{Reference trajectory}
\label{subsubsec:raceline}
Both runs follow the same optimised raceline (Fig.~\ref{fig:raceline}),
produced with the TUM global race trajectory optimisation
package~\cite{tumracetraj} under its minimum-curvature
formulation~\cite{heilmeier2020} on a recorded centreline of the CARLA
silverstone circuit: a closed \SI{5828.2}{\meter} lap of \num{2916} waypoints
at \SI{2.0}{\meter} spacing, reference speed
\SIrange{9.62}{25.78}{\meter\per\second}, minimum corner radius
\SI{18.7}{\meter} ($|\kappa|_{\max}=\SI{0.0536}{\per\meter}$). Its peak lateral
demand is $\max(v_{x}^{2}|\kappa|)=\SI{4.98}{\meter\per\second\squared}$
(\SI{0.51}{\g}), which saturates the flat
\SI{5.0}{\meter\per\second\squared} \texttt{ggv} envelope the optimiser was
given and stays at roughly \SI{73}{\percent} of the \SI{0.70}{\g} this
surface sustains. This
trajectory replaces the earlier one, whose \SI{1.22}{\g} demand the
admissibility gate of Section~\ref{sec:integration} rejected against the same
plant. The regeneration, not a relaxed gate, is what unblocked the closed-loop
runs below.

\begin{figure*}[t]
\centering
\includegraphics[width=\textwidth]{raceline.pdf}
\caption{The optimised raceline used for both runs, generated with the TUM
global race trajectory optimisation package~\cite{tumracetraj} under its
minimum-curvature formulation~\cite{heilmeier2020} from a recorded centreline
of the CARLA silverstone circuit. (a) Track layout, coloured by lateral demand, with the marked corner
($R=\SI{18.7}{\meter}$) the tightest on the lap.
(b) Reference speed profile, \SIrange{9.62}{25.78}{\meter\per\second}.
(c) Lateral demand $v_{x}^{2}|\kappa|$ along the lap against the flat
\SI{5.0}{\meter\per\second\squared} envelope handed to the optimiser (dashed)
and the grip the vehicle has on this surface (dotted,
$\approx\SI{0.70}{\g}$, measured). The
profile saturates the envelope and no segment exceeds the plant's grip.}
\label{fig:raceline}
\end{figure*}

\subsubsection{The two configurations}
\label{subsubsec:scenarios}
Two configurations of the same pipeline are run back to back on that
trajectory, three timed laps each, with an untimed out-lap from the spawn point
to the raceline's $s=0$. They differ in the five settings of
Table~\ref{tab:scenarios} and in nothing else: the vehicle, the map, the
trajectory, the controller, the data budget, the co-identification loop count
and the Pacejka box constraints are shared. \texttt{Baseline-NN-MSE} is the
inherited configuration (MLP residual, one-step MSE objective, local
trust-region fit, no friction warm start, undamped coefficient update)
evaluated at full scale on the
corrected rollout, so the comparison isolates the five settings rather than
reproducing the published pipeline in full (that comparison is
Section~\ref{subsec:baselinecmp}). \texttt{Ours} is the shipped configuration.
Because the five settings move together, the comparison below is a
configuration comparison and not a per-factor ablation;
Section~\ref{subsec:ablations} separates the settings that carry the result,
offline. It is also not a test
of this paper's central contribution, since the excitation-aware rollout of
Section~\ref{subsec:excitation} is active in \emph{both} arms, so what
Section~\ref{subsec:resultsclosedloop} measures is the five secondary settings
on top of that correction. The correction itself is evidenced in
Section~\ref{subsec:sweep}.

\begin{table}[t]
\caption{The two configurations.}
\label{tab:scenarios}
\centering
\begin{threeparttable}
\footnotesize
\setlength{\tabcolsep}{4pt}
\begin{tabular}{@{}lll@{}}
\toprule
Setting & \texttt{Baseline-NN-MSE} & \texttt{Ours}\\
\midrule
Residual architecture & MLP (\texttt{baseline}) & S4D (\texttt{s4})\\
Residual objective & One-step MSE & Physics-informed,
\eqref{eq:pinnloss}\\
Pacejka solver & Trust-region~\cite{branch1999trust} & Differential
evolution~\cite{storn1997de}\\
Friction warm start & Off & On; $D$ pinned per axle\\
Update relaxation $\beta$ & \num{1.00} (undamped) & \num{0.20}\\
\bottomrule
\end{tabular}
\begin{tablenotes}[flushleft]\scriptsize
\item Everything not listed is identical between the two, including the
excitation-aware rollout of Section~\ref{subsec:excitation}, which both use.
$\beta$ is the relaxation of \eqref{eq:relax}; \num{1.00} is the inherited
undamped update.
\end{tablenotes}
\end{threeparttable}
\end{table}

\subsubsection{Provenance of the tire prior}
\label{subsubsec:priorprovenance}
The prior $\phip^{0}$ that both arms start from, and that the Tikhonov term of
\eqref{eq:fitobj} pulls toward, is the measured tire of Table~\ref{tab:ref}. It
was obtained by fitting \eqref{eq:mf} to CARLA's own per-wheel telemetry
(\texttt{feedback/tire\_forces}: slip angle, lateral force, normal load) over a
dedicated limit-excitation run that reached $\alpha=\SI{0.37}{\radian}$, axle
$\mu=\num{1.05}$ and \SI{1.50}{\g}, by differential evolution, at
\SI{135}{\newton} front and \SI{85}{\newton} rear residual RMSE. Three
consequences have to be stated plainly. First, that is the same channel the
force and friction results below are scored against, so the prior is not
independent of the validation data. Second, the manoeuvre that produced it is
one the pipeline itself is not allowed to run. The raceline of
Section~\ref{subsubsec:raceline} peaks at \SI{0.51}{\g} and never approaches
the tire's peak, which is exactly the excitation deficit this paper is about.

Third, and this is what makes the runs below informative rather than
self-confirming, the prior was measured on a surface these runs do not use. It
carries $D=\num{1.009}/\num{1.002}$ and an axle slope of
\SI{13.2}{\kilo\newton\per\radian} front and
\SI{13.9}{\kilo\newton\per\radian} rear, against the
\num{0.70} and the \num{17.4} and \SI{16.6}{\kilo\newton\per\radian} this surface
realises (Section~\ref{subsec:resultstire}). Returning the prior unchanged is
therefore not a good outcome here: it is a \SI{44}{\percent} grip overstatement
and a \SIrange{16}{24}{\percent} stiffness deficit, and it would be visible in
every metric of Table~\ref{tab:comparison}. What the arm has to do to score
well is \emph{move}, in the right direction, from a prior that is wrong.

\subsubsection{What is measured}
Three families of quantity are recorded, each against the plant's own
telemetry: the identified tire model (coefficients, axle force, peak friction),
the one-step state prediction, and the closed-loop tracking of the MPC that
consumes the model. Every figure in this section is exported together with the
CSV holding the samples behind it, and the cross-run comparison reads those
CSVs rather than re-deriving any metric, so a figure and the table cannot
disagree.

\subsection{Identifiability of the peak factor against rollout excitation}
\label{subsec:sweep}

The identifiability claim of Section~\ref{subsec:excitation} is a claim about
the rollout configuration alone, so it is tested where the answer is known
exactly rather than in CARLA. A tire with a known peak factor $D$ is rolled out
through the synthetic steering sweep of stage~2, the shipped rollout with the
residual network set to zero, and fitted back with the shipped solver. The fit
is run \emph{unregularised} and started from the true coefficients, so nothing
it loses is lost to the prior, to the start point or to a local minimum.
Everything it loses is lost to excitation. Table~\ref{tab:sweep} and
Fig.~\ref{fig:sweep} report the grid.

Three axes are swept, rollout speed $v$, sweep amplitude
$\delta_{\mathrm{sweep}}$ and true $D$, and each cell is repeated over five
noise realisations. The two configuration axes are swept one at a time,
because $\mu_{\mathrm{reach}}$ is a function of both and a claim that it, and
not the speed, governs the fit cannot be tested by moving the speed alone. The
repetition is over noise rather than over the solver because the noise-free
fit is exactly seed-invariant here: five multi-start seeds return the same
coefficients to four decimal places, and so does a jitter width seven times
wider. That is the direct measurement behind the statement that the degeneracy
is not a solver problem. The injected noise
is Gaussian on the rolled-out states at a one-step persistence RMSE measured on
CARLA runs of this vehicle, $\sigma_{v_{y}}=\SI{0.009}{\meter\per\second}$ and
$\sigma_{\omega}=\SI{0.005}{\radian\per\second}$, the same order as the
persistence rows of Table~\ref{tab:comparison}, and is what the intervals in
Table~\ref{tab:sweep} are computed over. Because the front and rear fits of one
rollout share that rollout, the interval is taken over the five realisations
and not over the ten fits.

\begin{table*}[t]
\caption{Peak factor $\hat D$ recovered from a known tire, along both arguments
of $\mu_{\mathrm{reach}}$.}
\label{tab:sweep}
\centering
\begin{threeparttable}
\footnotesize
\setlength{\tabcolsep}{6pt}
\begin{tabular}{@{}rrcccc@{}}
\toprule
swept & $\mu_{\mathrm{reach}}$ & \multicolumn{4}{c}{$\hat D$ (railed of 8) at true $D$ of}\\
\cmidrule(l){3-6}
value & & 0.70 & 1.05 & 1.40 & 1.80\\
\midrule
\multicolumn{6}{@{}l}{\emph{(a) rollout speed $v$ [\si{\meter\per\second}] swept at $\delta_{\mathrm{sweep}}=\SI{0.40}{\radian}$}}\\
4 & 0.43 & \num{0.52}$\pm$\num{0.036} (1.8) & \num{0.57}$\pm$\num{0.013} (2.8) & \num{0.59}$\pm$\num{0.018} (3.8) & \num{0.50}$\pm$\num{0.001} (7)\\
6 & 0.96 & \num{0.68}$\pm$\num{0.001} (1) & \num{1.04}$\pm$\num{0.014} (1.2) & \num{1.22}$\pm$\num{0.066} (1.8) & \num{1.26}$\pm$\num{0.097} (2)\\
8 & 1.70 & \num{0.69}$\pm$\num{0.001} (1) & \num{1.02}$\pm$\num{0.001} (0.8) & \num{1.36}$\pm$\num{0.002} (0) & \num{1.79}$\pm$\num{0.025} (1.2)\\
10 & 2.66 & \num{0.70}$\pm$\num{0.000} (0) & \num{1.03}$\pm$\num{0.001} (0.4) & \num{1.37}$\pm$\num{0.002} (1) & \num{1.75}$\pm$\num{0.002} (0)\\
13 & 4.49 & \num{0.72}$\pm$\num{0.001} (0) & \num{1.05}$\pm$\num{0.002} (0) & \num{1.38}$\pm$\num{0.001} (0.4) & \num{1.77}$\pm$\num{0.002} (0.4)\\
15 & 5.98 & \num{0.76}$\pm$\num{0.009} (1) & \num{1.05}$\pm$\num{0.002} (0) & \num{1.39}$\pm$\num{0.002} (0) & \num{1.78}$\pm$\num{0.001} (0)\\
20 & 10.64 & \num{0.92}$\pm$\num{0.004} (1) & \num{1.14}$\pm$\num{0.005} (0.4) & \num{1.43}$\pm$\num{0.002} (0) & \num{1.80}$\pm$\num{0.005} (0)\\
\addlinespace
\multicolumn{6}{@{}l}{\emph{(b) $\delta_{\mathrm{sweep}}$ [\si{\radian}] swept at $v=\SI{15}{\meter\per\second}$}}\\
0.03 & 0.45 & \num{0.80}$\pm$\num{0.019} (2) & \num{1.10}$\pm$\num{0.041} (2) & \num{1.08}$\pm$\num{0.182} (2) & \num{1.09}$\pm$\num{0.397} (2)\\
0.06 & 0.90 & \num{0.92}$\pm$\num{0.003} (2) & \num{1.30}$\pm$\num{0.008} (2) & \num{1.68}$\pm$\num{0.017} (2) & \num{1.80}$\pm$\num{0.232} (3.4)\\
0.10 & 1.50 & \num{0.96}$\pm$\num{0.004} (2) & \num{1.34}$\pm$\num{0.016} (2) & \num{1.60}$\pm$\num{0.017} (1) & \num{1.97}$\pm$\num{0.030} (1.4)\\
0.15 & 2.24 & \num{0.99}$\pm$\num{0.004} (2) & \num{1.38}$\pm$\num{0.003} (2) & \num{1.71}$\pm$\num{0.019} (2) & \num{1.98}$\pm$\num{0.009} (1)\\
0.25 & 3.74 & \num{0.77}$\pm$\num{0.002} (2) & \num{1.12}$\pm$\num{0.004} (2) & \num{1.81}$\pm$\num{0.020} (2) & \num{2.00}$\pm$\num{0.000} (2)\\
0.40 & 5.98 & \num{0.76}$\pm$\num{0.009} (1) & \num{1.05}$\pm$\num{0.002} (0) & \num{1.39}$\pm$\num{0.002} (0) & \num{1.78}$\pm$\num{0.001} (0)\\
0.54 & 8.08 & \num{0.83}$\pm$\num{0.004} (1) & \num{1.08}$\pm$\num{0.002} (0) & \num{1.38}$\pm$\num{0.001} (0) & \num{1.78}$\pm$\num{0.008} (0)\\
\bottomrule
\end{tabular}
\begin{tablenotes}[flushleft]\scriptsize
\item Unregularised, started from the true coefficients. Each cell is $\hat D$ $\pm$ the half-width of its \SI{95}{\percent} $t$-interval over the five noise realisations, with the mean number of the eight coefficients resting on a box constraint in parentheses. The two axles of one rollout are not independent, so each realisation contributes the mean of its own $(\hat D_{f},\hat D_{r})$ and $n=5$; per-axle means and intervals are in the supplementary material.
\item $\mu_{\mathrm{reach}}$ is \eqref{eq:mureach} at $L=\SI{1.5334}{\meter}$; the inherited configuration is the first row of (a). It is the Ackermann \emph{upper} bound, and above $\mu_{\mathrm{reach}}\approx2$ it detaches from what the rollout realises: at $v=\SI{20}{\meter\per\second}$ it reads \num{10.64} while the rollout's realised $\max|v_{x}\omega|/g$ is \numrange{1.40}{2.18} across the four tires. A large $\mu_{\mathrm{reach}}$ means ``far past saturation'', not a demand the tire ever sees.
\end{tablenotes}
\end{threeparttable}
\end{table*}

\begin{figure*}[t]
\centering
\includegraphics[width=\textwidth]{identifiability_sweep.pdf}
\caption{The grid of Table~\ref{tab:sweep}, unregularised and started from the
true coefficients. (a) and (b) sweep the two arguments of
\eqref{eq:mureach} one at a time, one series per true $D$ (dotted lines), mean
$\pm$ \SI{95}{\percent} interval over the five noise realisations. Neither
axis is a speed effect: starving $\delta_{\mathrm{sweep}}$ at a fixed
\SI{15}{\meter\per\second} destroys $D$ as thoroughly as starving $v$, and the
\num{1.8} tire rails on its \num{2.0} bound while doing so. (c) reads both
against the demand the rollout \emph{realises},
$\max|v_{x}\omega|/g$, normalised by the tire's own $D$. Left of the dashed
vertical the rollout never reaches the peak and $\hat D$ reports the
excitation rather than the tire, low when the speed is the limiting factor
(filled) and high when the sweep amplitude is (open), and the two branches
close on
$\hat D/D=1$ from opposite sides as the realised demand passes $D$. The
box-edge counts are the parenthesised column of Table~\ref{tab:sweep}.}
\label{fig:sweep}
\end{figure*}

The grid says three things. First, the inherited
\SI{4}{\meter\per\second} rollout is uninformative about $D$ at \emph{every}
grip level tested. A $0.70$, a $1.05$, a $1.40$ and a $1.80$ tire all come
back as $\hat D$ between \num{0.50} and \num{0.59}, so the fit reports
$\mu_{\mathrm{reach}}$ and not the tire, with two to seven of the eight
coefficients on a bound. Second, the boundary is $\mu_{\mathrm{reach}}\gtrsim
D$ and not a speed: at \SI{6}{\meter\per\second}, where
$\mu_{\mathrm{reach}}=0.96$, the $0.70$ and $1.05$ tires are recovered to
within \SI{3}{\percent} while the $1.40$ and $1.80$ tires are
\SI{13}{\percent} and \SI{30}{\percent} short. That boundary is soft rather
than sharp. The $1.05$ tire is recovered to \SI{1}{\percent} at
$\mu_{\mathrm{reach}}=0.96$, slightly \emph{below} its own $D$, because a
rollout that stops just short of the peak still traverses the curvature
separating $B$, $C$ and $D$. The bound predicts where the information runs
out, not a step. Above $\mu_{\mathrm{reach}}=1.7$ every tire is recovered to
within \SI{4}{\percent} with no systematic railing. Third, the failure is not
the solver's. The noise-free fit is invariant to the multi-start seed and to
the jitter width, so multi-start, a global solver or a better start point
cannot recover a coefficient the data does not contain, which is why
Section~\ref{subsec:fit} reports the extra starts and the global solver as
secondary. The one caveat is at the top of the sweep, where a rollout far past
the peak begins to over-report a low-grip tire ($\hat D=\num{0.92}$ for a
$0.70$ tire at \SI{20}{\meter\per\second}), because the Magic Formula's
falling branch is being fitted to data the plant's own curve does not produce.
The extrapolation bound of Section~\ref{subsec:iterative} is what keeps the
shipped configuration away from that regime.

Panel (b) of Table~\ref{tab:sweep} sweeps the other argument of
\eqref{eq:mureach}: $\delta_{\mathrm{sweep}}$ from \SI{0.03}{\radian} to the
shipped $\delta_{\max}=\SI{0.54}{\radian}$, at a fixed
$v=\SI{15}{\meter\per\second}$. It confirms that speed is not the governing
variable, since starving the sweep amplitude destroys $D$ just as starving the
speed does, at the same fixed speed, and it corrects the shape of the claim in
two ways. First, the failure is not a common bias. Where an under-speeded rollout
returns \emph{too little} grip (every tire back at $\hat D\approx0.5$), an
under-swept one returns too much: at $\delta_{\mathrm{sweep}}=\SI{0.06}{\radian}$
the \num{1.05} tire comes back as \num{1.30} and the \num{1.80} tire rails on
the \num{2.0} bound, with two of the eight coefficients on a box edge in every
cell. What the two axes share is that $\hat D$ stops being a function of the
tire, the direction in which it then errs being set by the rollout: the
signature of an unidentified parameter rather than of a biased one.

Second, the two axes do \emph{not} collapse onto a common
$\mu_{\mathrm{reach}}$, and the reason is the upper-bound caveat of
Section~\ref{subsec:excitation}. Equation~\eqref{eq:mureach} uses the Ackermann
curvature, so it tracks the realised demand only while the vehicle's own
understeer is small; the realised $\max|v_{x}\omega|/g$ saturates near the
tire's peak while $\mu_{\mathrm{reach}}$ keeps growing, reaching \num{10.64}
against a realised \numrange{1.40}{2.18} at $v=\SI{20}{\meter\per\second}$.
Read against the realised demand (Fig.~\ref{fig:sweep}c) both axes do behave
alike: recovery sets in where the rollout actually reaches the tire's peak.
Equation~\eqref{eq:mureach} should therefore be used as the pipeline uses it,
as a cheap screening bound that is conclusive when it falls below $D$, and not
as a measure of the excitation achieved.

One property of this experiment bounds how all of it should be read. The
residual network is set to zero, so a Magic Formula is fitted to Magic-Formula
data and the model structure is exactly matched. The recovery reported here is
the best case, and the real pipeline's degradation at a given excitation is at
least this bad.

\subsection{Cross-run summary}
\label{subsec:summary}

Table~\ref{tab:comparison} is the complete comparison. \texttt{Ours} is better
on every quantity the benchmark scores against the plant (peak friction,
axle force on both axles, both one-step states, lateral tracking error in every
phase, and both error tails) and is behind only on the RMS QP solve time, by
\SI{5}{\percent} of a figure two orders of magnitude inside the control period.
The margins are large: \numrange{9.0}{12.1}$\times$ on peak friction and axle
force, \num{1.5}$\times$ on MPC tracking error. The rest of this section reads
those rows against the figures, and
Section~\ref{subsec:ablations} attributes them.

\begin{table}[t]
\caption{Cross-run comparison of the two configurations.}
\label{tab:comparison}
\centering
\begin{threeparttable}
\scriptsize
\setlength{\tabcolsep}{4pt}
\begin{tabular}{@{}lrr@{}}
\toprule
Metric & \texttt{Ours} & \texttt{Baseline-NN-MSE}\\
\midrule
\multicolumn{3}{@{}l}{\emph{Closed-loop tracking (3 timed laps)}}\\
Timed laps & 3 & 3\\
Run duration [\si{\second}] & 580.4 & 554.6\\
Best lap [\si{\second}], wall clock & 142.84 & 135.82\\
Mean lap [\si{\second}], wall clock & 154.47 & 145.39\\
RMS $e_{y}$, overall [\si{\meter}] & \textbf{0.256} & 0.319\\
RMS $e_{y}$, pure pursuit [\si{\meter}] & \textbf{0.387} & 0.443\\
RMS $e_{y}$, MPC [\si{\meter}] & \textbf{0.151} & 0.224\\
Max $|e_{y}|$ [\si{\meter}] & 2.628 & 2.628\\
Max $|e_{y}|$, MPC [\si{\meter}] & \textbf{0.683} & 0.776\\
ITAE $e_{y}$ [\si{\meter\second\squared}] & \textbf{16041} & 25502\\
IAE $e_{y}$ [\si{\meter\second}] & \textbf{60.2} & 87.2\\
IAE$/T$ [\si{\meter}] & \textbf{0.104} & 0.157\\
ITAE$/(T^{2}/2)$ [\si{\meter}] & \textbf{0.095} & 0.166\\
RMS heading error [\si{\radian}] & \textbf{0.0206} & 0.0215\\
Max $|$heading error$|$ [\si{\radian}] & 0.315 & \textbf{0.310}\\
Emergency halts & 0 & 0\\
MPC solve time, RMS [\si{\milli\second}] & 0.399 & \textbf{0.380}\\
MPC solve time, max [\si{\milli\second}] & 2.604 & \textbf{2.568}\\
\addlinespace
\multicolumn{3}{@{}l}{\emph{Axle lateral force against per-wheel telemetry}}\\
Front RMSE [\si{\newton}] & \textbf{27.2} & 329.6\\
Rear RMSE [\si{\newton}] & \textbf{33.2} & 348.4\\
Front $R^{2}$ & \textbf{0.993} & $-0.160$\\
Rear $R^{2}$ & \textbf{0.987} & $-0.461$\\
Front max abs.\ error [\si{\newton}] & \textbf{264} & 1295\\
Rear max abs.\ error [\si{\newton}] & \textbf{135} & 1316\\
\addlinespace
\multicolumn{3}{@{}l}{\emph{Peak friction against the plant's $\mu=0.70$}}\\
Front $\mu$ RMSE, whole run & \textbf{0.145} & 1.300\\
Rear $\mu$ RMSE, whole run & \textbf{0.142} & 1.300\\
Front / rear $\mu$ RMSE, after first release & \textbf{0.078} & 1.300\\
\addlinespace
\multicolumn{3}{@{}l}{\emph{One-step state prediction}}\\
$v_{y}$ RMSE [\si{\meter\per\second}] & \textbf{0.014} & 0.156\\
$\omega$ RMSE [\si{\radian\per\second}] & \textbf{0.049} & 0.104\\
$v_{y}$ persistence RMSE [\si{\meter\per\second}] & 0.011 & 0.017\\
$\omega$ persistence RMSE [\si{\radian\per\second}] & 0.005 & 0.007\\
Skill vs.\ persistence, $v_{y}$ & $\mathbf{-0.27}$ & $-8.18$\\
Skill vs.\ persistence, $\omega$ & $\mathbf{-8.8}$ & $-13.9$\\
\bottomrule
\end{tabular}
\begin{tablenotes}[flushleft]\scriptsize
\item Bold is the better of the two where the metric has a direction.
\item Lap times are on the ROS wall clock; the two runs did not hold the same real-time factor, so they are reported as recorded and not bolded.
\item Max $|e_{y}|$ is the same pure-pursuit out-lap transient in both runs, before either identified model is in the loop, and agrees to \SI{0.03}{\milli\meter}; the MPC-phase maximum beneath it is the one the identified model influences. ITAE is not scale-free and the runs differ in length, hence $\mathrm{IAE}/T$ and $\mathrm{ITAE}/(T^{2}/2)$ beside it.
\item The persistence rows are the reference baseline, next state equals last
measured state, and not a result.
\item ``After first release'' is the \SI{440}{\second} of each run following the first cycle on which the friction gate released a measured $\mu$; before it, \texttt{Ours} publishes the $\mu=\num{1.05}$ prior (Section~\ref{subsec:resultsmu}). The two axles carry the same $D$ by construction (Section~\ref{subsec:friction}), so their post-release errors coincide.
\item \textbf{Every value is a single run per configuration: no row carries a confidence interval, and the directional rows are one paired observation, not an estimate of a mean difference (Section~\ref{subsec:limits}).}
\end{tablenotes}
\end{threeparttable}
\end{table}

\subsection{Identified tire model}
\label{subsec:resultstire}

Figure~\ref{fig:pacejka} puts both identified curves against the plant's own
curve~\eqref{eq:physx} and lists the coefficients. Three things are visible.

First, no coefficient of \texttt{Ours} rests on a box constraint, in any of its
six accepted identifications. Its last accepted set is
$B=\num{12.890}/\num{13.746}$ inside $[4,20]$,
$C=\num{1.240}/\num{1.248}$ inside $[1.2,2.2]$, $D=\num{0.768}$ on both axles
inside $[0.4,2.0]$ and $E=\num{-2.517}/\num{-1.836}$ inside $[-3,1]$.
\texttt{Baseline-NN-MSE} puts $D$ on its upper bound of \num{2.000} on both
axles in \emph{every} accepted identification, and $E_{f}$ on its lower bound
of \num{-3.0} likewise; its first identification additionally rails $B$ on the
lower bound of \num{4.0} on both axles. A peak factor of \num{2.000} against a
plant that realises \num{0.70} is less a \SI{186}{\percent} error in a tire
parameter than the box-edge signature of
Section~\ref{subsec:identifiability}: the coefficient is not identified, and
nothing in a per-coefficient bounds check distinguishes that from a fit. The
one coefficient of \texttt{Ours} worth reading with care is $C$, which settles
\num{0.04} above its lower bound: inside the box, but close enough to it
that it should be reported as shaped by the constraint rather than as
identified, consistent with $C$'s role in the $BCD$ product of
Section~\ref{subsec:bcd}. $E$ should not be read as identified in either
configuration: it shapes only the region around and past the peak, which
neither the lap nor the bounded rollout reaches, and it is exactly the
coefficient the frozen Tikhonov target of \eqref{eq:fitobj} exists to hold near
its nominal value. It does so here, \texttt{Ours} ending at \num{-2.52}
against a \num{-2.00} prior, where the unregularised arm ends on the bound.

Second, and this is the change from the earlier configuration, the linear range
is now identified rather than inherited. Every slope here is
\eqref{eq:stiffness} for an identified set, and for the plant it is the slope
fitted to the simulator's own per-wheel forces over
$|\alpha|\in[\num{0.005},\num{0.02}]\,\si{\radian}$, all at the same static
axle loads $F_{z,f}=\SI{1373}{\newton}$, $F_{z,r}=\SI{1274}{\newton}$, so the
three are comparable. The plant realises
\SI{17.4}{\kilo\newton\per\radian} front and
\SI{16.6}{\kilo\newton\per\radian} rear over this run's load range, below the
$\bar F_{z}=1$ closed-form values of Table~\ref{tab:ref} for the reason its
footnote gives; the comparison here is against the realised slope, not against
the tabulated one. The configured prior
carries \num{13.2} and \num{13.9}, \SI{24}{\percent} and \SI{16}{\percent} low.
\texttt{Ours} returns \num{16.9} and \num{16.8}, i.e.\ \SI{3.2}{\percent} low
at the front and \SI{1.4}{\percent} high at the rear, so the identification
removes \SI{87}{\percent} of the prior's front stiffness error and
\SI{91}{\percent} of its rear one. The product
$BCDF_{z}$ that Section~\ref{subsec:bcd} identifies as the only quantity the
rollout fit determines is what moves and then holds: the first cycle takes the
front value from the prior's \SI{13.2}{\kilo\newton\per\radian} to
\num{14.6}, and the five cycles after the friction gate releases sit within
\SIrange{16.0}{16.9}{\kilo\newton\per\radian}, a spread of \SI{5}{\percent}
about a plant at \num{17.4}. \texttt{Baseline-NN-MSE} returns \num{37.9} and
\num{38.4}, \SI{118}{\percent} and \SI{132}{\percent} high, which is
$D=\num{2.0}$ multiplying a $B$ the fit chose against it.

Third, both sets keep the understeer sign ($l_{r}C_{r}>l_{f}C_{f}$) and so
pass the stability gate of Section~\ref{subsec:gates}, at margins of
\SI{7.3}{\percent} for \texttt{Ours} ($l_{f}C_{f}=\num{12.45}$ against
$l_{r}C_{r}=\SI{13.36}{\kilo\newton\meter\per\radian}$) and \SI{9.1}{\percent}
for \texttt{Baseline-NN-MSE} (\num{27.98} against \num{30.52}), against
\SI{2.5}{\percent} for the plant itself. The gate that does separate them is
the grip ceiling: \eqref{eq:gripceiling} gives \SI{2.00}{\g} for
\texttt{Baseline-NN-MSE} against \SI{0.77}{\g} for \texttt{Ours}, a plant
that sustains \SI{0.71}{\g}, and a trajectory that demands \SI{0.51}{\g}. The
baseline's ceiling is a claim the plant cannot honour by a factor of
\num{2.8}; the prior's \SI{1.01}{\g} is one it cannot honour by \num{1.4}; only
the identified set is admissible on the quantity the controller actually
consumes.

What makes all three readable on this run is that the lap does now traverse
part of the curvature. Over the force-validation samples,
$P_{99}(|\alpha|)=\SI{0.047}{\radian}$ (\SI{2.7}{\degree}) and
$\max|\alpha|=\SI{0.102}{\radian}$ (\SI{5.8}{\degree}) against the
\SI{8.0}{\degree} at which the plant's curve first reaches its peak at
$\mu=\num{0.70}$, and the axle utilisation $\max|F_{y}|/F_{z}$ reaches
\num{0.71} front and \num{0.72} rear. The same trajectory on the
$\mu=\num{1.05}$ surface of the earlier runs reached only \num{0.49}. The
excitation is still not a limit manoeuvre (the rollout of
Section~\ref{subsec:excitation} is what reaches up the curve, and the warm
start of Section~\ref{subsec:friction} is what carries independent information
about $\mu$), but it is enough for the friction gate of
Section~\ref{subsec:friction} to release, which on the $\mu=\num{1.05}$ runs
it did not.

\begin{figure}[t]
\centering
\includegraphics[width=\columnwidth]{pacejka_identified_vs_nominal_by_scenario}
\caption{Identified tire models against the plant. Top and middle: the plant's
own curve~\eqref{eq:physx} at the measured static load of each axle (black)
against each run's last accepted identification, over
$\alpha\in[-0.2,0.2]\,\si{\radian}$. Bottom: the identified Magic Formula
coefficients. \texttt{Baseline-NN-MSE}'s accepted peak factor sits on its
box constraint $D=\num{2.0}$ on both axles in every cycle, so its curve carries
roughly \num{2.8} times the plant's force at the edge of the plotted range,
while \texttt{Ours} tracks both the plant's initial slope and its peak. Note
that the per-wheel lateral force and the per-axle vertical load exported in the
underlying CSV are not on the same basis; the axle forces plotted here and
every force metric in Table~\ref{tab:comparison} are axle sums.}
\label{fig:pacejka}
\end{figure}

\subsection{Peak friction}
\label{subsec:resultsmu}

Table~\ref{tab:comparison} reports the error of the identified peak factor
against the effective friction the plant reports per wheel, and
Fig.~\ref{fig:mutimeseries} shows how it evolves over the run. This is the one
place where the pinned-$D$ design of Section~\ref{subsec:friction} is directly
observable, because the arm starts \SI{44}{\percent} away from the answer.

\texttt{Ours} publishes six accepted identifications. The first carries
$D_{f}=\num{1.009}$, $D_{r}=\num{1.002}$: the friction gate rejected that
buffer, so the axles held the prior they came in with, which is the
hold-rather-than-fit rule of Section~\ref{subsec:frictionfloor} and not a
fitted value. From $t=\SI{140.6}{\second}$ the gate releases on every
subsequent cycle and the published peak factor drops to
\num{0.779}, \num{0.821}, \num{0.749}, \num{0.756} and \num{0.768}, front and rear alike. The plant is at
\num{0.700}. The five released values therefore sit \SI{10.7}{\percent} high on
average with a spread of \num{0.072}, against a prior that was
\SI{44}{\percent} high. Over the whole run the $\mu$ RMSE is \num{0.145} front
and \num{0.142} rear, the first \SI{140}{\second} at the prior dominating it;
over the \SI{440}{\second} after the first release it is \num{0.078} on both
axles.

\texttt{Baseline-NN-MSE} returns $D=\num{2.000}$ on both axles in every
accepted identification and never leaves the bound, giving a constant
$\mu$ error of \num{1.300} on both axles. \texttt{Ours} is
\num{9.0}$\times$ better on the whole-run figure and \num{16.7}$\times$ better
once its warm start has released, and the difference is worth \SI{1.23}{\g} of
claimed grip against a plant that sustains \num{0.71}.

Two qualifications belong on the \texttt{Ours} column. The peak factor is
\emph{measured}, by the brush estimator of Section~\ref{subsec:friction} on the
recorded trace, and then held; it is not recovered by the Pacejka fit, which
Section~\ref{subsec:bcd} shows cannot recover it. And the residual
\SI{10.7}{\percent} bias is one-signed, consistent with the brush model's known
tendency to over-read $\mu$ where the curve is not fully traversed, which at
$\max|F_{y}|/F_{z}=\num{0.71}$ it is not.

\begin{figure*}[t]
\centering
\includegraphics[width=\textwidth]{mu_timeseries_by_scenario}
\caption{Peak friction over the run, one row per configuration: the plant's
effective $\mu$ (black) against the identified peak factor $D$ published at
each accepted identification. The surface is static, so every
departure from the black line is estimator error. \texttt{Ours} publishes the
$\mu=\num{1.05}$ prior while the friction gate is closed, then steps down onto
the plant within \SI{11}{\percent} at the first release and holds there for the
remaining five cycles, while \texttt{Baseline-NN-MSE} sits on the
$D\le\num{2.0}$ box constraint for the whole run.}
\label{fig:mutimeseries}
\end{figure*}

\subsection{Axle lateral force}
\label{subsec:resultsforce}

The identified curve is scored at the measured $(\alpha,F_{z})$ against the
plant's own lateral force, summed per axle. Table~\ref{tab:comparison} carries
the summary, and the bar charts, error histograms and time series behind it are
in the supplementary material. \texttt{Ours} leads on both axles and by an
order of magnitude: \SI{27.2}{\newton} against \SI{329.6}{\newton} RMSE at the
front and \SI{33.2}{\newton} against \SI{348.4}{\newton} at the rear, with
$R^{2}$ \num{0.993} and \num{0.987} against \num{-0.160} and \num{-0.461}. A
negative $R^{2}$ means the baseline's identified curve predicts the plant's
axle force worse than the mean of that force does, which is the force-metric
reading of a set whose initial slope is more than twice the plant's. The error
tails order the same way, maximum absolute error \SI{264}{\newton} and
\SI{135}{\newton} against \SI{1295}{\newton} and \SI{1316}{\newton}.

This is the linear-range result of Section~\ref{subsec:resultstire} seen as
force. At $P_{99}(|\alpha|)=\SI{2.7}{\degree}$ the force metric is dominated by
cornering stiffness and says comparatively little about the peak, so the two
models separate here almost entirely on the slope, and
\texttt{Baseline-NN-MSE}'s slope carries its railed $D=\num{2.0}$. The peak
factor still has consequences here: because $D$ multiplies the slope
in \eqref{eq:stiffness}, a $D$ on its bound is the direct cause of the force
error rather than an overstatement of grip the linear range hides, which is the
opposite of what the earlier $\mu=\num{1.05}$ runs suggested.

The comparison is also the reason the relaxed update of \eqref{eq:relax}
belongs in the method rather than in a configuration file. Pinning $D$ alone
does not produce the numbers above; with $\beta=1$ the same arm degrades
monotonically with iteration count, and the offline sweep of
Section~\ref{subsec:ablations} separates the two.

\subsection{One-step state prediction}
\label{subsec:resultsstate}

Level~1 of the validation protocol (Section~\ref{subsec:validation}) is
reported in Table~\ref{tab:comparison} against a persistence baseline, next
state equals last measured state, written by the same benchmark node.
\texttt{Ours} leads on both signals, lateral velocity at
\SI{0.014}{\meter\per\second} RMSE against \SI{0.156}{\meter\per\second}
($R^{2}$ \num{0.998} against \num{0.709}) and yaw rate at
\SI{0.049}{\radian\per\second} against \SI{0.104}{\radian\per\second}
($R^{2}$ \num{0.894} against \num{0.514}).

Neither model beats persistence at the \SI{0.035}{\second} model step. The
skill scores $1-\mathrm{RMSE}/\mathrm{RMSE}_{\mathrm{pers}}$ are negative for
both configurations and both signals, worst at \num{-13.9}
(\texttt{Baseline-NN-MSE}, $\omega$) and best at \num{-0.27}
(\texttt{Ours}, $v_{y}$). We report this rather than the $R^{2}$
alone because $R^{2}$ against a slowly varying signal is close to
uninformative. A one-step metric on a smooth trajectory rewards not moving,
and over one \SI{0.035}{\second} step the plant's own states barely move
($\mathrm{RMSE}_{\mathrm{pers}}$ is \SI{0.011}{\meter\per\second} and
\SI{0.005}{\radian\per\second}). One-step accuracy is consequently weak
evidence for either tire model, and the gap here is large enough to be read
only as confirming the direction the axle-force and closed-loop comparisons
establish, not as an independent result.
The two state traces against their own ground truth over the whole run are in
the supplementary material.

\subsection{Closed-loop tracking}
\label{subsec:resultsclosedloop}

Level~3 of the validation protocol completes on the regenerated trajectory.
Both runs finish three timed laps with zero emergency halts and the same
finite-state sequence (bootstrap pure pursuit, running pure pursuit, handover,
running MPC), with one transition of each kind and no fallback to pure pursuit
after handover. The sequence is read off each run's exported state timeline,
which records the handover at $t=\SI{79.1}{\second}$ for \texttt{Ours} and
$t=\SI{64.0}{\second}$ for \texttt{Baseline-NN-MSE} and no state change
thereafter. The aggregated \texttt{transitions} column of the cross-run
summary is not populated and is omitted from Table~\ref{tab:comparison}. Under MPC, which is where the identified model enters the
loop, \texttt{Ours} tracks at \SI{0.151}{\meter} RMS lateral error against
\SI{0.224}{\meter}, a \SI{32}{\percent} reduction, with maximum
\SI{0.683}{\meter} against \SI{0.776}{\meter}
(Fig.~\ref{fig:trackingbox}). Overall ITAE is \num{16041} against \num{25502}.
ITAE is $\int t|e_{y}|\,\mathrm{d}t$ and is neither
scale-free nor comparable across runs of different length, and these two runs
are \SI{580.4}{\second} and \SI{554.6}{\second} long, so it is reported here
beside two normalised forms computed from the same samples: mean absolute
error $\mathrm{IAE}/T$, \SI{0.104}{\meter} against \SI{0.157}{\meter}, and
$\mathrm{ITAE}/(T^{2}/2)$, \SI{0.095}{\meter} against \SI{0.166}{\meter}. All
three order the runs the same way and at the same magnitude, so the
normalisation does not carry the result. The ITAE split confirms that the
whole difference sits in the MPC phase: the pure-pursuit phases, which do not
use the identified model, give \num{110} and \num{107} respectively, within
\SI{3}{\percent} of each other and two orders of magnitude smaller. The signed
error distribution (supplementary material) shows the mechanism as a narrower
distribution about a comparable centre, at sample standard deviation
\SI{0.256}{\meter} against \SI{0.318}{\meter} and mean \SI{0.021}{\meter}
against \SI{0.024}{\meter}.

The QP cost is unchanged. RMS solve time is \SI{0.399}{\milli\second} against
\SI{0.380}{\milli\second} and worst case \SI{2.60}{\milli\second} against
\SI{2.57}{\milli\second}, both two orders of magnitude below the
\SI{50}{\milli\second} control period, and both slightly in the baseline's
favour at a margin of \SI{5}{\percent} and \SI{1}{\percent}. The identified
coefficients change the model the QP carries, not its size.

Lap time is not read as a result. Timing runs on the ROS wall clock while the
simulator advances its fixed step faster than real time, and the ratio between
the two differed between the runs, so the recorded lap times in
Table~\ref{tab:comparison} are not simulated lap times and their
\SIrange{5}{6}{\percent} spread is not a vehicle-performance difference.

\begin{figure}[t]
\centering
\includegraphics[width=\columnwidth]{tracking_error_by_scenario}
\caption{Absolute lateral tracking error by configuration and active
controller. The pure-pursuit boxes cover the acquisition phase before handover
and do not use the identified model. The MPC boxes do.}
\label{fig:trackingbox}
\end{figure}

\subsection{Baseline pipeline versus this work}
\label{subsec:baselinecmp}

The supplementary material tabulates the pipeline as
published~\cite{ontrack2025} against the pipeline as it stands here, on every
axis that differs. In summary, the residual model, the data budget, the
four stages and the steady-state force recovery are inherited unchanged.
What differs is the configuration of the virtual-data step, the relaxation on
the loop closing it, the regularisation
topology, the measurement chain, and the contract at the output.
Table~\ref{tab:scenarios}'s comparison is a subset of those axes: the rollout
correction of Section~\ref{subsec:excitation} is active in both runs of
Section~\ref{subsec:summary}, so that section measures the remaining five
settings and not the rollout.

\subsection{Peak-friction warm start on a synthetic plant}
\label{subsec:mu}

The friction warm start of Section~\ref{subsec:friction} is evaluated on a
synthetic single-track plant with brush tires rather than in CARLA, because
only a synthetic plant supplies a $\mu$ that is known exactly and an excitation
level that can be swept independently of it. The plant is driven by a
swept-sine steering input for \SI{40}{\second} at \SI{50}{\hertz} with true
$\mu=1.05$ and $C_{\alpha}=12.1F_{z}$, and the steering amplitude sets how far
up the tire curve the manoeuvre reaches. Table~\ref{tab:mu} reports what each
estimator returns.

\begin{table}[t]
\caption{Peak friction recovered from a synthetic brush-tire plant (true $\mu=1.05$), against excitation.}
\label{tab:mu}
\centering
\begin{threeparttable}
\small
\begin{tabular}{@{}cccc@{}}
\toprule
$\delta$ amplitude & Grip utilisation & Utilisation & Brush fit
\eqref{eq:brush}\\
\midrule
\SI{0.05}{\radian} & \SI{45}{\percent}  & 0.473 & \textbf{1.048}\\
\SI{0.07}{\radian} & \SI{63}{\percent}  & 0.658 & \textbf{1.048}\\
\SI{0.09}{\radian} & \SI{80}{\percent}  & 0.834 & \textbf{1.048}\\
\SI{0.13}{\radian} & \SI{100}{\percent} & 1.045 & \textbf{1.048}\\
\bottomrule
\end{tabular}
\begin{tablenotes}[flushleft]\scriptsize
\item The utilisation estimator tracks the grip the manoeuvre \emph{used},
while the brush fit of \eqref{eq:brush} recovers the peak from the curvature
of $F_{y}(\alpha)$ and is therefore flat in excitation.
\item Below \SI{40}{\percent} utilisation the brush fit is gated off and the utilisation bound is reported instead.
\end{tablenotes}
\end{threeparttable}
\end{table}

At \SI{45}{\percent} utilisation, a level an ordinary lap reaches, the
utilisation bound is \SI{55}{\percent} low while the brush fit is within
\SI{0.2}{\percent}. Across true $\mu\in\{0.7,1.05,1.4\}$ the fit stays within
\SI{5}{\percent} and $C_{\alpha}$ within \SI{10}{\percent}, and with odometry
and IMU noise injected $\mu$ stays within \SI{12}{\percent}. Below the
\SI{40}{\percent} floor the gate suppresses the fit and reports the bound, so
the estimator's failure is visible rather than silent. On the CARLA run of
Section~\ref{subsec:resultsmu} the gate behaves the same way at both ends: it
rejects the first buffer, and on the five it accepts it returns $\mu$ within
\SI{11}{\percent} of the plant at an absolute axle utilisation of
\numrange{0.54}{0.68}. The \SI{11}{\percent} residual bias against the
\SI{0.2}{\percent} of the synthetic plant is the cost of a real rollout, a
reconstructed $I_{z}\dot\omega$ term and a curve the manoeuvre traverses only
to \num{0.71} of $F_{z}$, and it is the dominant term in the identified peak
factor's error.

\subsection{What carries the result}
\label{subsec:ablations}

The two arms of Section~\ref{subsec:summary} differ in five settings at once,
so that comparison establishes an ordering and attributes nothing. This
subsection attributes it, at the cost of leaving the loop: every number below
is an \emph{offline replay} of the identification stage on recorded buffers,
not a closed-loop run.

\subsubsection{Protocol}
\texttt{nn\_train} is replayed on the five \SI{30}{\second} identification
buffers each arm collected during its own run of
Section~\ref{subsec:summary}, and each configuration is scored with the
quantity the closed-loop tire-force benchmark itself reports: the identified
Magic Formula evaluated at the plant's own measured slip angles, against the
plant's own measured lateral forces, summed per axle. No synthetic plant and no
re-derived metric are involved, so the numbers here and in
Table~\ref{tab:comparison} are on the same scale. One caveat applies to anyone
reproducing this from the exported curve-validation files: the lateral force
they carry is \emph{per wheel} while the vertical load beside it is the
\emph{axle} static load, and fitting them as written returns every cornering
stiffness a factor of two low.

\subsubsection{Damping the co-identification loop}
Table~\ref{tab:beta} sweeps $\beta$ of \eqref{eq:relax} at six iterations
throughout. Two columns are the argument. At $\beta=\num{1.00}$, the inherited
undamped update, axle force RMSE is \num{70.4} and \SI{62.3}{\newton}, and it is
reached by degrading monotonically from \num{38.9} and \SI{26.3}{\newton} at a single
iteration: the loop is not refining the estimate, it is amplifying its own
extrapolation. At $\beta=\num{0.00}$ the coefficients never move from the
pinned prior, which is the pipeline performing no identification at all, and
that is the \emph{worst} column of the sweep at
\num{92.5} and \SI{71.8}{\newton}. Between the two, the minimum sits at
$\beta=\num{0.20}$ (\SI{23.5}{\newton} $\pm$ \SI{1.5}{\newton} over seeds
\num{0}, \num{3} and \num{7}), and
the identified $BCD$ rises monotonically with $\beta$ across the whole sweep,
so the damping is trading how much of each fit is accepted against how much
extrapolation error comes with it, rather than simply interpolating toward the
prior. Six iterations are therefore defensible on evidence, but only damped.

\begin{table}[t]
\caption{Relaxation sweep. Five recorded buffers, six iterations, offline replay.}
\label{tab:beta}
\centering
\begin{threeparttable}
\scriptsize
\setlength{\tabcolsep}{3.2pt}
\begin{tabular}{@{}lrrrrrrrrr@{}}
\toprule
$\beta$ & 0.00 & 0.05 & 0.10 & 0.15 & \textbf{0.20} & 0.25 & 0.30 & 0.50 & 1.00\\
\midrule
Front $F_{y}$ RMSE [\si{\newton}] & 92.5 & 68.3 & 48.6 & 33.9 & \textbf{28.0} & 31.2 & 36.4 & 55.6 & 70.4\\
Rear $F_{y}$ RMSE [\si{\newton}] & 71.8 & 49.0 & 31.3 & 19.8 & \textbf{19.4} & 26.0 & 30.2 & 38.4 & 62.3\\
Front $BCD$ [\si{\per\radian}] & 7.82 & 8.82 & 9.66 & 10.36 & 10.88 & 11.32 & 11.65 & 12.75 & 12.42\\
\bottomrule
\end{tabular}
\begin{tablenotes}[flushleft]\scriptsize
\item $\beta=\num{1.00}$ is the inherited undamped update and $\beta=\num{0.00}$ the pinned prior, i.e.\ no identification. \num{0.20} ships.
\end{tablenotes}
\end{threeparttable}
\end{table}

\subsubsection{Pinning the peak factor}
Replayed on five buffers of a $\mu=\num{1.05}$ run, letting the fit move $D$
gives a mean $|D-\mu|$ of \num{0.585} and a maximum of \num{0.950}, with the
identified front stiffness spread over \SI{32.9}{\kilo\newton\per\radian}
across the five. Pinning only the axles whose brush fit was released reduces
those to \num{0.191}, \num{0.378} and \num{32.3}; additionally holding the
axles with no released measurement, which is the shipped rule, reduces them to
\num{0.084}, \num{0.213} and \SI{18.5}{\kilo\newton\per\radian}. The middle row
is the informative one: releasing an axle helps only where the gate fires, and
what removes most of the remaining error is refusing to let an unmeasured $D$
be fitted at all. The same replay is the source of the $BCD$ degeneracy
measurement quoted in Section~\ref{subsec:bcd}.

\subsubsection{Normalising the physics-informed objective}
The unnormalised form of \eqref{eq:pinnloss} puts \SI{99.9996}{\percent} of the
objective on the dynamics-residual term at initialisation, against
\SI{0.0004}{\percent} on the data; with the $T_{s}/m$ and $T_{s}/I_{z}$ scaling
of \eqref{eq:physnorm} the shares become \SI{1.47}{\percent} and
\SI{98.53}{\percent}. The correction is necessary but not sufficient on its
own: applied without the damping it makes the force metric slightly
\emph{worse} (\num{77.1} and \SI{76.2}{\newton} against
\num{70.4} and \SI{62.3}{\newton}), because an objective that now fits the data also
supplies a sharper trajectory for the undamped loop to amplify. With the
damping in place, the two arrive at
\num{28.0} and \SI{19.4}{\newton} and $R^{2}=\num{0.986}$ at the front, against
\num{92.5} and \SI{71.8}{\newton} and \num{0.870} for the unidentified prior.

\subsubsection{Two settings that do less than claimed}
Two of the five settings of Table~\ref{tab:scenarios} deserve a weaker
statement than the closed-loop comparison alone would support. The S4D residual
of Section~\ref{subsec:s4} is worth \SI{21}{\percent} on the force metric at
matched settings and a single iteration, \SI{35.3}{\newton} against the
baseline MLP's \SI{44.7}{\newton} as a mean over the two axles. That is the
ablation the earlier version of this work listed as unrun, and its sign is
positive, but it is an offline replay on recorded buffers and not a
closed-loop result. The residual \emph{objective} is a wash: with the
normalisation and the damping both in place, \texttt{physics\_informed} and
plain \texttt{mse} are indistinguishable at \SI{23.7}{\newton} each. We keep
the physics-informed objective because it costs nothing and is the paper's
stated formulation, and we do not claim it buys accuracy on this manoeuvre,
because it measurably does not.

\subsection{Identification cost}
\label{subsec:runtime}

One S4D identification, six co-identification iterations on a \SI{30}{\second}
lap, takes \SI{20.9}{\second} end to end, inside the \SI{30}{\second}
re-identification period the node runs at. The reduction from the
\SI{181.4}{\second} it originally took is exact: no epoch count, iteration
count, early-stopping rule or loss weight was changed, and the identified
coefficients agree to four decimal places. The per-phase breakdown,
the device split, and the caveat that the stage micro-benchmarks and the
end-to-end totals are separate measurements rather than a decomposition of each
other, are in the supplementary material.
